\section{Task 4: Graph Neural Network}

\begin{questionbox}
\textbf{Q1.} Do you still have to pad your data? If yes, how?
\end{questionbox}

No. In the GNN representation, each sample is a graph with its own number of nodes and edges.
Mini-batches are built as disjoint unions of graphs, so there is no need to pad all samples to the
same number of API calls, nodes, or edges.

\begin{questionbox}
\textbf{Q2.} Do you have to truncate the testing sequences? Justify your answer.
\end{questionbox}

No. GNN layers apply the same message-passing functions to all nodes and edges, independently of the
graph size. Therefore, test sequences longer than the training ones can be converted into larger
graphs and processed without truncation, as long as unknown API calls are mapped to \texttt{<UNK>}.

\begin{questionbox}
\textbf{Q3.} What is the advantage of modeling your problem with a GNN with respect to an RNN in
this scenario? However, what do you lose?
\end{questionbox}

A GNN models the transition structure of the API trace directly. Nodes represent API calls and
directed edges represent consecutive transitions, including loops when they occur. This avoids
sequential unrolling and can capture local transition patterns more explicitly than an RNN.

However, the graph is a static, unweighted representation. We lose the exact temporal position of
each call and repeated transition counts: two traces can induce the same unweighted graph even if
they execute transitions in different orders or with different frequencies.

\begin{table}[h!]
\centering
\begin{tabular}{lcccccc}
\hline
\textbf{Split} & \textbf{Samples} & \textbf{Avg. nodes} & \textbf{Avg. edges} &
\textbf{Avg. loops} & \textbf{Max nodes} & \textbf{Max edges} \\
\hline
Train & 16325 & 21.947 & 39.643 & 0.000 & 41 & 72 \\
Test  & 6505  & 24.279 & 45.104 & 0.000 & 45 & 76 \\
\hline
\end{tabular}
\caption{Statistics of the transition graphs induced by train and test API sequences.}
\label{tab:task4_graph_stats}
\end{table}

\begin{questionbox}
\textbf{Q4.} Train and tune variations of GNNs considering different message and aggregation
functions and architectures: Simple GCN, GraphSAGE, and GAT. How does each model perform with
respect to the previous architectures? Can you beat the baseline?
\end{questionbox}

We trained three graph models: a simple GCN, GraphSAGE, and GAT. API IDs are mapped to learnable node
embeddings, then processed with message-passing layers. Hyperparameters were selected with Optuna on
a stratified validation split from the training set, optimizing macro F1. The search space is
reported in Table~\ref{tab:task4_optuna_space}.

\begin{table}[h!]
\centering
\begin{tabular}{llll}
\hline
\textbf{Hyperparameter} & \textbf{Search space} & \textbf{Scale/type} & \textbf{Models} \\
\hline
\texttt{embedding\_dim} & \{16, 32, 64\} & Categorical & All \\
\texttt{hidden\_dim} & \{64, 128, 256\} & Categorical & All \\
\texttt{num\_layers} & \{2, 3\} & Categorical & All \\
\texttt{dropout} & $[0, 0.5]$ & Uniform & All \\
\texttt{batch\_size} & \{64, 128, 256\} & Categorical & All \\
\texttt{lr} & $[10^{-5}, 3 \cdot 10^{-3}]$ & Log-uniform & All \\
\texttt{weight\_decay} & $[10^{-6}, 10^{-2}]$ & Log-uniform & All \\
\texttt{aggregation} & \{\texttt{mean}, \texttt{max}\} & Categorical & GraphSAGE \\
\texttt{heads} & \{2, 4\} & Categorical & GAT \\
\hline
\end{tabular}
\caption{Optuna search space for the GNN models.}
\label{tab:task4_optuna_space}
\end{table}

\begin{table}[h!]
\centering
\begin{tabular}{lcccccccc}
\hline
\textbf{Model} & \textbf{Emb.} & \textbf{Hidden} & \textbf{Layers} & \textbf{Dropout} &
\textbf{LR} & \textbf{Weight decay} & \textbf{Batch} & \textbf{Extra} \\
\hline
GCN       & 32 & 64  & 3 & 0.3006 & 0.000568 & 0.000001 & 64  & -- \\
GraphSAGE & 32 & 256 & 2 & 0.0028 & 0.001047 & 0.000672 & 128 & mean \\
GAT       & 32 & 256 & 2 & 0.0028 & 0.001047 & 0.000672 & 128 & 2 heads \\
\hline
\end{tabular}
\caption{Best hyperparameters found for the three GNN variants.}
\label{tab:task4_best_hyperparameters}
\end{table}

\begin{figure}[h!]
\centering
\begin{minipage}{0.32\textwidth}
\centering
\includegraphics[width=\textwidth]{figures/task-4/gcn_training_curve.pdf}
\end{minipage}
\hfill
\begin{minipage}{0.32\textwidth}
\centering
\includegraphics[width=\textwidth]{figures/task-4/sage_training_curve.pdf}
\end{minipage}
\hfill
\begin{minipage}{0.32\textwidth}
\centering
\includegraphics[width=\textwidth]{figures/task-4/gat_training_curve.pdf}
\end{minipage}
\caption{Validation macro F1 curves for GCN (left), GraphSAGE (center), and GAT (right).}
\label{fig:task4_training_curves}
\end{figure}

\begin{figure}[h!]
\centering
\includegraphics[width=0.95\textwidth]{figures/task-4/loss_curves.pdf}
\caption{Training and validation loss curves for the three GNN variants.}
\label{fig:task4_loss_curves}
\end{figure}

\begin{table}[h!]
\centering
\begin{tabular}{lcccccc}
\hline
\textbf{Model} & \textbf{Best val F1} & \textbf{Accuracy} & \textbf{Balanced acc.} &
\textbf{Macro F1} & \textbf{F1 Goodware} & \textbf{F1 Malware} \\
\hline
GCN       & 0.7384 & 0.9491 & 0.8094 & 0.7324 & 0.4916 & 0.9732 \\
GraphSAGE & 0.8244 & 0.9702 & 0.8302 & 0.8071 & 0.6298 & 0.9845 \\
GAT       & 0.8253 & 0.9736 & 0.7667 & 0.7959 & 0.6055 & 0.9863 \\
\hline
\end{tabular}
\caption{Final test performance of the three GNN variants.}
\label{tab:task4_test_performance}
\end{table}

\begin{table}[h!]
\centering
\begin{tabular}{lcccc}
\hline
\textbf{Model/Class} & \textbf{Precision} & \textbf{Recall} & \textbf{F1-score} & \textbf{Support} \\
\hline
GCN Goodware (0)       & 0.3922 & 0.6584 & 0.4916 & 243 \\
GCN Malware (1)        & 0.9864 & 0.9604 & 0.9732 & 6262 \\
GraphSAGE Goodware (0) & 0.5872 & 0.6790 & 0.6298 & 243 \\
GraphSAGE Malware (1)  & 0.9875 & 0.9815 & 0.9845 & 6262 \\
GAT Goodware (0)       & 0.6839 & 0.5432 & 0.6055 & 243 \\
GAT Malware (1)        & 0.9824 & 0.9903 & 0.9863 & 6262 \\
\hline
\end{tabular}
\caption{Per-class classification report for the GNN variants.}
\label{tab:task4_classification_report}
\end{table}

\begin{figure}[h!]
\centering
\begin{minipage}{0.32\textwidth}
\centering
\includegraphics[width=\textwidth]{figures/task-4/gcn_confusion_matrix.pdf}
\end{minipage}
\hfill
\begin{minipage}{0.32\textwidth}
\centering
\includegraphics[width=\textwidth]{figures/task-4/sage_confusion_matrix.pdf}
\end{minipage}
\hfill
\begin{minipage}{0.32\textwidth}
\centering
\includegraphics[width=\textwidth]{figures/task-4/gat_confusion_matrix.pdf}
\end{minipage}
\caption{Confusion matrices for GCN (left), GraphSAGE (center), and GAT (right).}
\label{fig:task4_confusion_matrices}
\end{figure}

All GNN variants beat the frequency-based Logistic Regression baseline in macro F1. The best GNN was
GraphSAGE, with macro F1 0.8071 and goodware F1 0.6298. GAT achieved the highest accuracy and malware
F1, but its lower goodware recall reduced balanced accuracy and macro F1. GCN was the weakest GNN,
although it still improved over the baseline macro F1.

\begin{table}[h!]
\centering
\begin{tabular}{lcccc}
\hline
\textbf{Model} & \textbf{Accuracy} & \textbf{Balanced acc.} & \textbf{Macro F1} & \textbf{Goodware F1} \\
\hline
Logistic Regression baseline & 0.9688 & 0.6654 & 0.7154 & 0.4469 \\
Best FFNN: Embedding FFNN    & 0.9356 & 0.5789 & 0.5749 & 0.1832 \\
Best RNN: BiLSTM             & 0.9748 & 0.8663 & 0.8381 & 0.6894 \\
Best GNN: GraphSAGE          & 0.9702 & 0.8302 & 0.8071 & 0.6298 \\
\hline
\end{tabular}
\caption{Comparison between the best GNN and the previous architectures on the test set.}
\label{tab:task4_architecture_comparison}
\end{table}

Compared with the previous architectures, GraphSAGE is much better than the FFNN models and better
than the Logistic Regression baseline, but it does not beat the best RNN. The BiLSTM remains the best
overall model in macro F1 and goodware F1, suggesting that the exact temporal order still carries
useful information that is partly lost when the sequence is collapsed into an unweighted transition
graph.
